\section{Frullani's Integral}

\begin{enumerate}

    % --- Theorem statement and basic conditions ---

  \item State Frullani's integral theorem.
  \item What conditions on $f(x)$ are typically required for Frullani's theorem to hold?
  \item Why must the limits $f(0)$ and $f(\infty)$ exist?
  \item What is the role of the condition $a,b>0$?
  \item Why does the logarithm naturally appear in Frullani-type integrals?
  \item Derive
    \[
      \int_0^\infty \frac{e^{-ax} - e^{-bx}}{x}\,dx = \ln\frac{b}{a}.
    \]
  \item Explain why
    \[
      \int_0^\infty \frac{e^{-ax}}{x}\,dx
    \]
    diverges.
  \item Why does subtracting two individually divergent integrals sometimes produce a convergent result?

    % --- Edge cases ---

  \item What happens if $a = b$?
  \item What happens if one of the parameters tends to zero?
  \item What happens if one of the parameters tends to infinity?

    % --- Differentiation under the integral sign (Feynman trick) ---

  \item Show how Frullani's theorem can be derived using differentiation under the integral sign.
  \item Define
    \[
      F(a) = \int_0^\infty \frac{e^{-x} - e^{-ax}}{x}\,dx.
    \]
    Compute $F'(a)$, use the initial condition $F(1) = 0$ to determine the constant of integration, and write
    down the closed form for $F(a)$.
  \item When applying the Feynman trick to
    \[
      I(a) = \int_0^\infty \frac{f(x) - f(ax)}{x}\,dx,
    \]
    what initial condition pins down the constant of integration? Why does this condition hold?
  \item Differentiate
    \[
      \int_0^\infty \frac{f(ax) - f(bx)}{x}\,dx
    \]
    with respect to $a$, solve the resulting equation, and verify that the answer matches Frullani's theorem.
  \item In the Feynman approach, differentiating with respect to $a$ removes the $\tfrac{1}{x}$ factor from
    the integrand. Why does this make the integral tractable, and what condition on $f$ justifies
    differentiating under the integral sign?
  \item Why does the final answer depend only on $f(0)$ and $f(\infty)$ rather than the full shape of $f$?
  \item Give an example of a function for which Frullani's theorem fails, and identify which hypothesis is violated.

    % --- Alternative derivations and computations ---

  \item Derive Frullani's theorem from the identity
    \[
      f(ax) - f(bx) = \int_b^a x\,f'(tx)\,dt.
    \]
  \item Compute
    \[
      \int_0^\infty \frac{\arctan(ax) - \arctan(bx)}{x}\,dx.
    \]
  \item Compute
    \[
      \int_0^\infty \frac{\ln(1+ax) - \ln(1+bx)}{x}\,dx.
    \]
    Verify the answer by checking the $a = b$ limit.
  \item Does Frullani's theorem apply to $\displaystyle\int_0^\infty \frac{\cos(ax) - \cos(bx)}{x}\,dx$?
    Identify which hypothesis is (or is not) satisfied, and compute the integral by another method.
  \item What changes if the integration interval is $(0,T)$ rather than $(0,\infty)$?

    % --- Laplace transform connection ---

  \item Write $\dfrac{1}{x} = \int_0^\infty e^{-xs}\,ds$ and substitute into $\displaystyle\int_0^\infty
    \frac{f(ax)-f(bx)}{x}\,dx$.
    After a formal interchange of the order of integration, express the result in terms of the Laplace
    transform $\mathcal{L}\{f\}(s) = \int_0^\infty f(x)e^{-sx}\,dx$.
  \item Show that
    \[
      \int_0^\infty \frac{f(ax)-f(bx)}{x}\,dx
      = \int_0^\infty \left[
        \frac{1}{a}\,\mathcal{L}\{f\}\!\left(\frac{s}{a}\right)
        - \frac{1}{b}\,\mathcal{L}\{f\}\!\left(\frac{s}{b}\right)
      \right]ds.
    \]
    Verify this for $f(x) = e^{-x}$ by computing the right-hand side directly and confirming the answer
    equals $\ln(b/a)$.
  \item State the identity
    \[
      \int_0^\infty \frac{g(x)}{x}\,dx = \int_0^\infty \mathcal{L}\{g\}(s)\,ds
    \]
    and give sufficient conditions for its validity. Use it to rederive Frullani's theorem in one step
    by setting $g(x) = f(ax) - f(bx)$ and computing $\mathcal{L}\{g\}(s)$ explicitly.
  \item In the Laplace-transform approach, why is the condition $g(0) = 0$ necessary for the identity
    in the previous question? Show that $g(0) = 0$ follows automatically when $g(x) = f(ax) - f(bx)$ and
    the hypotheses of Frullani's theorem are satisfied.
  \item What justifies the interchange of the order of integration in the Laplace-transform proof? State
    the theorem you invoke (e.g.\ Tonelli's theorem) and verify its hypotheses for $f(x) = e^{-x}$.
  \item How is Frullani's theorem connected to Laplace transforms in a broader sense? (\textit{Hint}:
      consider the formula $\int_0^\infty f(x)/x\,dx = \int_0^\infty F(s)\,ds$, where $F = \mathcal{L}\{f\}$,
    and how it applies once the condition $g(0) = 0$ is ensured.)

    % --- Comparison: Feynman trick vs.\ Laplace transform ---

  \item Both the Feynman trick and the Laplace-kernel method deal with the $\tfrac{1}{x}$ factor. In the
    Feynman approach, $\tfrac{1}{x}$ is absorbed by differentiating with respect to a parameter; in the
    Laplace approach, $\tfrac{1}{x}$ is represented as $\int_0^\infty e^{-xs}\,ds$. Explain the structural
    analogy between these two approaches, and describe a setting in which one is preferable to the other.

    % --- Assumptions ---

  \item What assumptions are required to differentiate under the integral sign? State a sufficient condition
    (e.g.\ dominated convergence) and verify it for $f(x) = e^{-x}$ in the Feynman proof of Frullani's theorem.

\end{enumerate}
